% =========================================================
% 1.1 PLANTEAMIENTO DEL PROBLEMA
% El \section{} de esta sección está en el chapter.tex del capítulo;
% sus referencias van en seccion1_1.bib (misma carpeta).
% =========================================================

La inseguridad constituye una de las principales problemáticas urbanas de la Ciudad de México, afectando de manera particular la movilidad peatonal de sus habitantes. De acuerdo con la Encuesta Nacional de Seguridad Pública Urbana (ENSU) de junio de 2026, el $69\%$ de la población de la alcaldía Gustavo A. Madero declaró sentirse insegura, cifra que supera en casi diez puntos porcentuales la media nacional de $59.8\%$ \cite[p. 5]{2026ensu}.

Esta percepción varía dependiendo de las características del lugar y de la cantidad de personas que lo utilizan. Como se muestra en el Cuadro \ref{tab:percepcion_inseguridad}, basado en la Encuesta Nacional de Victimización y Percepción sobre Seguridad Pública (ENVIPE) 2025 y la ENSU 2026, los espacios de tránsito cotidiano como los cajeros automáticos en la vía pública, el transporte público y las calles presentan los mayores niveles de percepción de inseguridad a nivel nacional \cite[p. 31]{ENVIPE2025}; \cite[p. 8]{2026ensu}.


\begin{table}[H]
    \centering
    \caption[Percepción de inseguridad ciudadana por tipo de espacio]{Percepción de inseguridad ciudadana por tipo de espacio a nivel nacional \cite{ENVIPE2025, 2026ensu}.}
    \vspace{0.2cm}
    \begin{tabular}{|l|c|c|c|}
        \hline
        \textbf{Tipo de Espacio} & \textbf{2025} & \textbf{Marzo 2026} & \textbf{Junio 2026} \\ \hline
        Cajero automático en vía pública & 73.5\% & 70.6\% & 70.3\% \\ \hline
        Transporte público & 65.1\% & 64.1\% & 62.4\% \\ \hline
        Calle (vía pública abierta) & 62.0\% & 65.3\% & 63.4\% \\ \hline
        Parque o centro recreativo & 49.9\% & 46.7\% & 45.4\% \\ \hline
        Centro comercial & 39.7\% & 34.4\% & 32.1\% \\ \hline
        Trabajo & 25.8\% & 28.2\% & 26.9\% \\ \hline
        Casa & 16.2\% & 16.1\% & 15.8\% \\ \hline
        Escuela & 33.6\% & 15.2\% & 14.3\% \\ \hline
    \end{tabular}
    \label{tab:percepcion_inseguridad}
\end{table}


Estos datos muestran que los espacios públicos forman parte importante de la percepción de inseguridad de la población y es por esta razón que el análisis del riesgo peatonal no puede basarse únicamente en los delitos que son registrados por las autoridades. También es importante considerar cómo perciben
las personas la seguridad de los lugares por los que transitan. 

En los años 2017-2024, el robo o asalto en la calle y el transporte público registró una tasa de 6,003 delitos por cada 100 mil habitantes \cite[p. 11]{ENVIPE2025}. Sin embargo, las cifras oficiales no muestran todos los delitos que ocurren. De acuerdo con las encuestas de victimización, el 90.4 \% de los delitos ocurridos no fueron denunciados. Al considerar también aquellos casos en los que hubo una denuncia pero no se inició una carpeta de investigación, la llamada cifra negra alcanzó el 93.2 \% en 2025 \cite[p. 23]{ENVIPE2025}. Esta situación también se presenta en los delitos que ocurren en la vía pública. Por ejemplo, la ENVIPE 2022 estimó una cifra Negra cercana al 94 \% para el robo a transeúnte a nivel nacional \cite{alvaro_Obregon_cifra}..

La falta de denuncias puede deberse a diferentes factores, como la poca confianza en las autoridades, el tiempo que implica realizar una denuncia o la percepción de que recuperar los bienes robados es poco probable. Por esta razón, los registros oficiales deben entenderse como una parte del problema y no como una representación completa de los delitos que ocurren en el espacio público.

A pesar de esta problemática, los peatones no cuentan actualmente con herramientas tecnológicas que les permitan evaluar de manera objetiva el nivel de riesgo asociado a una ruta antes de recorrerla. Las aplicaciones de navegación que se usan actualmente usan exclusivamente variables de distancia y tiempo de traslado, sin incorporar ningún criterio relacionado con la seguridad del entorno urbano. Por otra parte, los sistemas existentes que sí intentan representar el riesgo delictivo lo hacen fundamentalmente a partir de estadísticas oficiales de incidencia delictiva, es decir, de los delitos que efectivamente fueron denunciados.

En consecuencia, cualquier modelo que dependa únicamente de carpetas de investigación como fuente de riesgo estará ignorando sistemáticamente a la gran mayoría de los eventos delictivos realmente ocurridos, así como a las condiciones del entorno físico que, sin haber derivado en un delito denunciado, pueden representar igualmente una amenaza para el peatón.

Por todo lo anterior, se identifica la necesidad de desarrollar una herramienta que permita a los peatones de zonas de alta afluencia y vulnerabilidad, como es el caso de las Zonas Territoriales 6, 7 y 8 de la alcaldía Gustavo A. Madero, generar rutas que consideren de manera simultánea el tiempo de traslado y las condiciones de riesgo del entorno urbano, superando así la dependencia exclusiva de estadísticas delictivas oficiales.
